% ============================================
% Experimental Materials and Protocols
% ============================================
\subsection*{A. Experimental Materials and Protocols}
\label{sec:suppl:materials}

Full experimental protocols for bacterial strain isolation, media preparation, community assembly, and coalescence experiments are provided in the main text Methods section. Here we provide supplementary details on strain characteristics.

% --------------------------------------------
% A.1 Bacterial Strain Library
% --------------------------------------------
\subsubsection*{A.1 Bacterial Strain Library}

The 54 bacterial isolates used in this study were derived from soil, tree surface, and flower stamen environments collected in Cambridge, MA, USA. Phylogenetic identities were determined by full-length 16S rRNA gene sequencing (Supplementary Table~S1). The isolates spanned 29 families across three phyla: Proteobacteria, Firmicutes, and Bacteroidota (Supplementary Fig.~S1). Each isolate was purified by serial streaking on agar plates before storage as glycerol stocks (25\% glycerol, $-80$~°C).

% --------------------------------------------
% A.2 Strain Phenotypic Diversity
% --------------------------------------------
\subsubsection*{A.2 Strain Phenotypic Diversity}

Monoculture growth characteristics were measured for all 54 isolates prior to community assembly experiments. Growth yield ranged from OD$_{600}$ = 0.1 to 1.2 after 24~h, and final pH in unbuffered medium ranged from 4.5 to 8.5, indicating substantial phenotypic diversity in growth rate and environmental modification potential (Supplementary Fig.~S30).

